\section{Threats to Validity}
\label{sec:threats}

\subsection{Internal Validity}

% \textit{Alternative explanation of the reverse failure.}
% A fine-tuned model may infer a forward task from any Java input~\cite{kotha2024implicit}.
% It would then ignore the reverse instruction.
% Three observations limit this threat.
% First, the reverse outputs keep the size and structure of the obfuscated input.
% A model that re-applied the forward transformation would insert more dead code or encrypt more literals.
% Second, the reasoning trace on Identifier Renaming follows the reverse instruction.
% Third, held-out tasks lose less than the inverse with the same checkpoint.
% These observations do not fully separate the two explanations.
% Our definition of the effect is behavioral.
% Both explanations describe a model that keeps the trained direction under a reverse prompt.

\textit{Alternative explanation of the reverse failure.}
The reverse failure may not reflect a lost capability.
After fine-tuning, the model may simply infer the forward task from any Java input and ignore the reverse instruction~\cite{kotha2024implicit}.
Under this explanation, the failure would concern instruction following rather than the inverse itself.
Three observations argue against it.
First, a model that re-applied the forward transformation would change the obfuscated input further, for instance by inserting more dead code or encrypting more literals.
The reverse outputs show no such change, since they keep the size and structure of the obfuscated input (Section~\ref{sec:reverse-eval}).
Second, on Identifier Renaming, the reasoning trace follows the reverse instruction and proposes the correct mapping, even though the final output leaves the code unchanged (Section~\ref{sec:cs_mechanism}).
The model therefore reads the reverse request, but does not act on it.
Third, with the same checkpoint, unrelated tasks lose far less than the inverse, which suggests that instruction following is largely preserved.
These observations weaken the alternative explanation.
%A control that gives other Java tasks to the fine-tuned models would separate the two accounts.
This limit does not affect our main claim, because we define the inversion curse by behavior rather than by cause.
Both accounts describe a model that stays in the trained direction under a reverse prompt, so the alternative explanation is a possible cause of the effect.

\textit{Training configuration.}
The reverse failure could result from overly aggressive training rather than from forward fine-tuning itself.
We limit this risk in two ways.
First, we selected the training size and schedule through a preliminary study, and we kept the configuration with the best forward performance.
Second, we used standard LoRA settings for 7B models~\cite{hu2022lora} and the provider settings for the API models.
We did not sweep learning rates or LoRA ranks, and a milder configuration might reduce the loss.
The reverse failure nevertheless appears in both training regimes, with 1,000 pairs for the API models and 10,000 pairs for the open-source models.
These regimes also differ in model family and optimizer, so this comparison does not isolate the effect of training size.

\textit{Prompt sensitivity.}
Other prompts might elicit the inverse.
We tested four prompting strategies.
They include few-shot examples and chain-of-thought.
None restored the reverse direction after fine-tuning.

\textit{Data contamination.}
% Pretraining data may contain the evaluation programs.
% We evaluate on ConDefects~\cite{condefects}. % [EXISTING]
% This dataset was built to reduce leakage concerns.
% Contamination would also affect the baseline and the fine-tuned model alike.
% It cannot by itself produce a drop after fine-tuning.
% \textit{Data contamination.}
To limit contamination between adaptation and evaluation, we fine-tune on CodeNet and evaluate on ConDefects dataset.
Both datasets may nevertheless appear in the pretraining data of the models.
We reduce this risk on the fine-tuning side by generating the training pairs with an obfuscator.
The obfuscated programs are produced for this study, so their exact form is unlikely to occur in pretraining data.
Our results also vary with the transformation type and the training strategy, which points to an effect of adaptation on obfuscated data rather than to recall of memorized programs.

\subsection{Construct Validity}

% \textit{Reverse recovery criterion.}
% Our criterion counts any positive movement toward the source on any channel.
% This choice accepts many valid deobfuscations (Definition~\ref{def:pair}).
% It may also count small changes as recoveries.
% The baseline reverse rates may therefore overstate full recovery.
% This leniency does not weaken the post-fine-tuning result.
% Under the same lenient criterion, no fine-tuned configuration recovers any program.
% We also report the per-channel similarity values.
% Readers can thus judge the size of each change.

\textit{Reverse recovery criterion.}
Evaluating deobfuscation is a challenge.
Tests can confirm that an output preserves the behavior of the program, but the obfuscated input already passes the same tests.
Deciding whether an output is actually deobfuscated therefore requires syntactic measures, and to our knowledge no single metric establishes this in a systematic way.
We therefore use four complementary channels, each capturing a different aspect of recovery.
An output counts as a recovery if it is valid and moves closer to the original, and we read the channels together to judge how consistent a change is.
This broad rule accepts the many valid forms of a deobfuscated program, but it also counts small, partial changes.

\textit{Semantic equivalence.}
We use test suites as the oracle of equivalence.
Tests cannot cover every behavior.
Each evaluation program comes with an average of 16 tests.
Test-based equivalence is standard in code generation and translation studies~\cite{pan2024lost}.

\textit{Similarity and readability metrics.}
CodeBLEU, AST overlap, and readability metrics are proxies for recovery.
As mention before, no single metric captures deobfuscation quality.
We therefore combine several independent channels.
The readability model was built from human judgments~\cite{scalabrino2018readability}. 

\textit{Reasoning traces.}
Chain-of-thought text may not reflect the computation behind an answer~\cite{turpin2023unfaithful,jacovi2020faithfulness}.
We use traces as behavioral evidence only.
We make no claim about internal representations.

\subsection{External Validity}

\textit{Language and transformations.}
The main study uses Java programs and three obfuscation types from one tool.
These types cover lexical, structural, and data-flow transformations, and they match techniques observed in real Android apps~\cite{dong2018android, wermke2018obfuscation}.
They do not cover all existing techniques.
Industrial obfuscators also apply transformations such as control-flow flattening and reflection~\cite{dong2018android}, which we leave for future work.
Our programs are also file-level, and they do not fully reflect real-world code such as Android applications or multi-file repositories.
Such systems involve dependencies across files and often exceed the context size of current LLMs, which prevented us from evaluating them.
The translation and compilation checks extend the observation to two other paired tasks, but each check uses a single model, so their results are indicative.

\textit{Models.}
\label{sec:models-ttv}
We acknowledge that our study does not include current state-of-the-art models. Two constraints explain this choice.
First, fine-tuning access to proprietary models is limited.
%Anthropic does not offer it on its own platform~\cite{aws2024haikufinetuning}.
%Google restricts it to selected Gemini versions~\cite{google2026geminituning}.
OpenAI offers it for a few models only~\cite{azure2026finetuning}.

Second, recent open-weight models are far larger than the models we fine-tuned.
Table~\ref{tab:open-sizes} lists examples.
Only a fraction of their weights is active for each token.
Fine-tuning still requires the full set of weights in memory.
% At 16-bit precision, each parameter takes two bytes.
% The smallest model in Table~\ref{tab:open-sizes} thus needs about 460\,GB for its weights alone.
% A 7B model needs about 14\,GB.
LoRA reduces the number of trainable parameters~\cite{hu2022lora},
%QLoRA reduces memory through quantization~\cite{dettmers2023qlora}. % [EXISTING]
but does not remove the need to load the base model.
%This requirement exceeds our academic hardware. 
The smallest model in Table~\ref{tab:open-sizes} thus needs about 460\,GB for its weights alone.

\begin{table}[t]
\centering
\caption{Size of recent open-weight models. Memory is the approximate size of the weights at 16-bit precision (2 bytes per parameter).}
\label{tab:open-sizes}
\small
\begin{tabular}{lrrr}
\toprule
Model & Total & Active & Weights (16-bit) \\
\midrule
DeepSeek-V4-Pro~\cite{deepseek2026v4} & 1.6T & 49B & $\approx$3.2\,TB \\
Kimi K2~\cite{kimi2025k2} & 1T & 32B & $\approx$2.0\,TB \\
GLM-5~\cite{glm2026glm5} & 744B & 40B & $\approx$1.5\,TB \\
Mistral Large 3~\cite{mistral2025large3} & 675B & 41B & $\approx$1.35\,TB \\
Qwen3.5-397B-A17B~\cite{qwen2026qwen35} & 397B & 17B & $\approx$0.8\,TB \\
MiniMax-M2~\cite{minimax2026m2} & 230B & 10B & $\approx$460\,GB \\
\midrule
7B model (this study) & 7B & 7B & $\approx$14\,GB \\
\bottomrule
\end{tabular}
\end{table}

Our results may therefore not directly generalize to these models.
However, the effect appears in both open-source evaluated models and proprietary models, far larger (GPT series).
This suggests that the effect is not specific to open 7B models under LoRA.
%It does not show that the effect holds for current frontier models.

We also do not evaluate agentic systems.
Such systems use a model together with tools, whereas fine-tuning adapts the model itself, which is the object of our study.
An agent built on a model that has lost the inverse may inherit this limitation, although tools may compensate for part of it.

\textit{Adaptation methods.}
We study LoRA and provider fine-tuning.
Other adaptation methods exist, such as full fine-tuning and continual-learning methods like replay or EWC~\cite{kirkpatrick2017overcoming, weyssow2023continual}.
We therefore do not claim that our findings generalize to all these methods.

